\subsection{Integer Factorization}\label{sec:trapdoor_functions:integer_factorization}

The \textbf{\gls{if} problem} is a trapdoor function where the ``easy part'' is computing multiplications between prime numbers while the ``hard part'' is finding which prime numbers were multiplied to obtain a given number (i.e., factoring --- see \Cref{sec:mathematical_underpinnings:factorization}). More formally, given a large composite number $n$ --- where $n$ is the product of (usually) two prime numbers $p$ and $q$, i.e., $n = p \times q$ --- is it (believed to be) hard to find two (or more) prime factors of $n$.

The \gls{if} problem is the trapdoor function underlying:
\begin{itemize}
	\item the \gls{rsa} cipher (\Cref{sec:asymmetric_cryptography:ciphers:rsa}).
\end{itemize}

\exampleBlock{\gls{if} example}{It is easy to compute $7 \times 13 = 91$. However, given $91$, it is difficult to understand that the factorization of $91 = 7 \times 13$. In fact, the simplest approach to factorize a number is through trial divisions, where we try to divide the number with increasingly greater prime numbers hoping to find a factor. Of course, doing trial divisions for 91 is quick, but doing trial divisions for numbers with hundreds of digits is not.}

\warningBlock{The \gls{if} problem vs. quantum computing}{Quantum computing can solve the \gls{if} problem using the Shor's algorithm \cite{shor_algorithms_1994} in polynomial time (that is, efficiently).}
%TODO --- see \Cref{future_sec:post_quantum_cryptography}.}
